\chapter{Tavole Sinottiche e Comparazione degli Algoritmi}
\label{ch:appendice-quadro-sinottico-algoritmi}

\FloatBarrier
\section{Guida alla Consultazione e Modello di Riferimento}
\label{sec:guida-consultazione-sinottico}

La presente appendice raccoglie e sistematizza l'insieme delle metodologie algoritmiche, dei paradigmi di calcolo e delle strutture dati analizzati nel corso di \textit{Algorithm Engineering} attraverso cinque quadri sinottici tematici:
\begin{enumerate}
  \item \textbf{Ordinamento e Selezione:} algoritmi interni, esterni (DAM model) e basati su campionamento (Tabella~\ref{tab:sinottico-ordinamento-selezione});
  \item \textbf{Dizionari, Strutture Randomizzate e Campionamento:} alberi $B$-tree, skip lists, treap, alberi cartesiani e campionamento casuale uniforme con $n$ noto (Algorithm S, single-item) o incognito (Priority sampling, Reservoir R e Z) (Tabella~\ref{tab:sinottico-dizionari-streaming});
  \item \textbf{Intersezione di Posting Lists e Partizionamento:} approcci merge, dicotomici, partizionamento mutuale e skipping a blocchi a due livelli (Tabella~\ref{tab:sinottico-intersezione-liste});
  \item \textbf{Ordinamento e Strutture per Stringhe in RAM:} multi-key quicksort, radix sort, TST e trie compatti (Tabella~\ref{tab:sinottico-stringhe-ram});
  \item \textbf{Memoria Gerarchica, Indicizzazione e Sottostringhe:} blind search, suffix tree, suffix array e algoritmo di Kasai (Tabella~\ref{tab:sinottico-indicizzazione-sottostringhe}).
\end{enumerate} 

Per ciascun algoritmo vengono formalizzate:
\begin{itemize}
  \item \textbf{Complessit\`a Temporale:} valutata nel caso medio ($\mathbb{E}[T]$) e nel caso pessimo ($T_{\text{worst}}$), considerando sia il numero di confronti sia le operazioni CPU.
  \item \textbf{Complessit\`a Spaziale:} memoria primaria (RAM) ausiliaria strettamente necessaria all'esecuzione, oltre allo spazio per contenere i dati.
  \item \textbf{Complessit\`a di I/O (DAM Model):} numero di trasferimenti a blocchi tra memoria secondaria e memoria primaria secondo il Disk Access Machine Model di Aggarwal e Vitter~\cite{aggarwal1988}, ove:
  \begin{equation}
    \text{Scan}(N) = \Theta\left(\frac{N}{B}\right), \qquad \text{Sort}(N) = \Theta\left(\frac{N}{B} \log_{M/B} \frac{N}{M}\right)
  \end{equation}
  con $N$ taglia totale del dataset, $M$ capacit\`a della RAM e $B$ dimensione del blocco fisico di trasferimento.
  \item \textbf{Note Operative e Trade-off:} contesto applicativo ideale, presupposti architetturali e colli di bottiglia hardware.
\end{itemize}

\FloatBarrier
\section{Ordinamento e Selezione in Memoria Primaria e Gerarchica}
\label{sec:sinottico-ordinamento-selezione}

La Tabella~\ref{tab:sinottico-ordinamento-selezione} sintetizza gli algoritmi di ordinamento esterno e selezione discussi nei Capitoli da~\ref{cha:lecture-three} a~\ref{chap:external-distribution-sort}.

\begin{table}[H]
\centering
\small
\resizebox{\linewidth}{!}{%
\begin{tabular}{>{\raggedright\arraybackslash\bfseries}p{4.2cm} c >{\raggedright\arraybackslash}p{3.2cm} >{\raggedright\arraybackslash}p{2.5cm} >{\raggedright\arraybackslash}p{3.4cm} >{\raggedright\arraybackslash}p{6.6cm}}
\toprule
\textbf{Algoritmo / Tecnica} & \textbf{Cap.} & \textbf{Tempo CPU} & \textbf{Spazio RAM} & \textbf{I/O Disco (DAM)} & \textbf{Note Operative e Punti di Forza} \\
\midrule
2-Way External Merge Sort & \ref{cha:lecture-three} & $\mathcal{O}(N \log N)$ & $\mathcal{O}(B)$ & $\mathcal{O}\left(\frac{N}{B} \log_2 \frac{N}{M}\right)$ & Paradigma Divide-et-Impera base: genera run da $M$ elementi e li fonde a coppie. Subottimale con RAM ampia. \\
\midrule
Multiway Merge Sort ($d$-way) & \ref{cha:lecture-four} & $\mathcal{O}(N \log N)$ & $\mathcal{O}(M)$ & $\mathcal{O}\left(\frac{N}{B} \log_{M/B} \frac{N}{M}\right)$ & \textbf{Ottimo asintotico $\text{Sort}(N)$}: sfrutta un ventaglio $d \approx M/B - 1$, minimizzando i passi di fusione ad albero. \\
\midrule
Replacement Selection (Snow Plow) & \ref{cha:lecture-four} & $\mathcal{O}(N \log(M/B))$ & $\mathcal{O}(M)$ (Min-Heap) & Invariato ($\text{Sort}(N)$) & Algoritmo dello spazzaneve: produce run di lunghezza media attesa pari a $\mathbf{2M}$, dimezzando i run iniziali. \\
\midrule
Multi-Disk Striping (RAID-0) & \ref{cha:lecture-four} & $\mathcal{O}(N \log N)$ & $\mathcal{O}(D \cdot B)$ & $\mathcal{O}\left(\frac{N}{D \cdot B} \log_{M/B} \frac{N}{M}\right)$ & Distribuisce i blocchi in round-robin su $D$ dischi indipendenti, scalando il throughput di I/O di un fattore $D$. \\
\midrule
Quicksort Dual-Pivot (Yaroslavskiy) & \ref{cha:lecture-five} & $\approx 1.9 n \ln n$ / $\mathcal{O}(n^2)$ & $\mathcal{O}(\log n)$ stack & Cache-friendly & Due pivot e 3 partizioni; riduce i cache miss L1/L2 e il traffico di memoria rispetto a Quicksort classico. \\
\midrule
Bounded Quicksort & \ref{cha:lecture-five} & $\mathcal{O}(N \log N)$ & $\mathcal{O}(M)$ & $\mathcal{O}\left(\frac{N}{B} \log_2 \frac{N}{M}\right)$ & Arresta la ricorsione esterna quando i sottoproblemi hanno dimensione $\le M$, ordinandoli in RAM con 0 I/O. \\
\midrule
Bounded Max-Heap (Top-$k$) & \ref{cha:lecture-six} & $\mathcal{O}(n \log k)$ & $\mathcal{O}(k)$ heap & $0$ I/O (in RAM) & Mantiene un Max-Heap di $k$ elementi; eccellente per $k \ll n$, evitando l'ordinamento globale del dataset. \\
\midrule
Quickselect (Hoare) & \ref{cha:lecture-six} & $\mathcal{O}(n)$ medio / $\mathcal{O}(n^2)$ & $\mathcal{O}(\log n)$ ricors. & $\mathcal{O}(N/B)$ streaming & Scende nel solo ramo del pivot che racchiude il rango target; media teorica di $3.38\,n$ confronti atomici. \\
\midrule
Quickselect con Campionamento & \ref{cha:lecture-six} & $\mathcal{O}(n)$ medio & $\mathcal{O}(\log n)$ & $\mathcal{O}(N/B)$ streaming & Seleziona il pivot su un piccolo campione casuale, abbattendo i confronti attesi a $\approx 2n + o(n)$. \\
\midrule
Median-of-Medians (BFPRT) & \ref{cha:lecture-six} & $\mathcal{O}(n)$ deterministico & $\mathcal{O}(\log n)$ & $\mathcal{O}(N/B)$ & Partiziona a gruppi di 5; garantisce sbilanciamento massimo $30/70$. Costanti elevate per uso pratico. \\
\midrule
External Distribution Sort & \ref{chap:external-distribution-sort} & $\mathcal{O}(N \log N)$ & $\mathcal{O}(M)$ & $\mathcal{O}\left(\frac{N}{B} \log_{M/B} \frac{N}{M}\right)$ & Duale del Merge Sort: campiona $a \cdot k$ pivot con over-sampling $a = \Theta(\ln k)$ garantendo bucket bilanciati con alta probabilit\`a. \\
\bottomrule
\end{tabular}%
}
\caption{Quadro comparativo degli algoritmi di ordinamento e selezione in memoria primaria e gerarchica.}
\label{tab:sinottico-ordinamento-selezione}
\end{table}

\FloatBarrier
\section{Dizionari, Strutture Dati Randomizzate e Flussi di Dati}
\label{sec:sinottico-dizionari-streaming}

La Tabella~\ref{tab:sinottico-dizionari-streaming} confronta le strutture dinamiche per dizionari e gli algoritmi di campionamento casuale uniforme su flussi di dati, distinguendo chiaramente tra gli scenari in cui la taglia della popolazione $n$ \`e nota a priori (accesso casuale o streaming a singolo passaggio con Algorithm S) oppure incognita (Priority sampling, Reservoir sampling classico e a salti) (Capitoli~\ref{cha:lecture-two}, \ref{chap:external-distribution-sort}, \ref{chap:random-sampling-data-streams}, \ref{chap:randomized-data-structures}, \ref{chap:randomized-search-trees-treaps}).

\begin{table}[H]
\centering
\small
\resizebox{\linewidth}{!}{%
\begin{tabular}{>{\raggedright\arraybackslash\bfseries}p{4.2cm} c >{\raggedright\arraybackslash}p{3.2cm} >{\raggedright\arraybackslash}p{2.5cm} >{\raggedright\arraybackslash}p{3.4cm} >{\raggedright\arraybackslash}p{6.6cm}}
\toprule
\textbf{Struttura / Algoritmo} & \textbf{Cap.} & \textbf{Tempo Operativo} & \textbf{Spazio RAM} & \textbf{I/O Disco (DAM)} & \textbf{Note Operative e Punti di Forza} \\
\midrule
$B$-Tree / $B^+$-Tree & \ref{cha:lecture-two} & $\mathcal{O}(\log_B N)$ per op. & Nodi su disco & $\mathcal{O}(\log_B N)$ & Albero bilanciato a grado elevato proporzionato a $B$; garantisce saturazione $\ge 50\%$ e range query ottime su foglie collegate ($B^+$). \\
\midrule
Skip List (Pugh) & \ref{chap:randomized-data-structures} & Search/Insert: $\mathcal{O}(\log n)$ att. / $\mathcal{O}(n)$ w.c. & $\mathcal{O}(n)$ nodi attesi & RAM-oriented & Liste collegate a pi\`u livelli generate tramite lanci di moneta; evita rotazioni complesse offrendo parallelismo lock-free. \\
\midrule
Treap (Aragon-Seidel) & \ref{chap:randomized-search-trees-treaps} & Search/Insert: $\mathcal{O}(\log n)$ att. / $\mathcal{O}(n)$ w.c. & $\mathcal{O}(n)$ nodi & RAM-oriented & Combina ordinamento BST sulle chiavi con priorit\`a Max-Heap random; supporta \textsc{Split} e \textsc{Merge} in $\mathcal{O}(\log n)$ atteso. \\
\midrule
Albero Cartesiano (Stack Monotono) & \ref{chap:randomized-search-trees-treaps} & Costruzione: $\mathcal{O}(n)$ deterministica & $\mathcal{O}(n)$ & RAM-oriented & Costruito linearmente con stack; trasforma query RMQ su intervalli continui in query di Lowest Common Ancestor (LCA) $\mathcal{O}(1)$. \\
\midrule
Rejection Sampling ($n$ noto) & \ref{chap:external-distribution-sort} & $\mathcal{O}(m)$ atteso ($m \ll n$) & $\mathcal{O}(m)$ dizionario & Accesso casuale RAM & Estrae indici uniformi in $[1, n]$ rigettando duplicati; presuppone $n$ noto e array indirizzabile ad accesso casuale. Degrada per $m \to n$; inapplicabile su flussi continui. \\
\midrule
Single-Item Sampling ($m=1$, $n$ noto) & \ref{chap:random-sampling-data-streams} & $\mathcal{O}(n)$ w.c. / $n/2$ atteso & $\mathcal{O}(1)$ ausiliario & $1$ passaggio streaming & \textbf{Streaming con $n$ noto ($m = 1$)}: seleziona l'elemento $j$-esimo se $p \le \frac{1}{n - j + 1}$ e arresta subito l'esecuzione. Assicura probabilit\`a uniforme esatta $\mathbb{P} = 1/n$ senza memorizzare il flusso. \\
\midrule
Algorithm S (Knuth, $n$ noto) & \ref{chap:random-sampling-data-streams} & $\mathcal{O}(n)$ tempo, $n$ chiamate RNG & $\mathcal{O}(1)$ ausiliario & $1$ passaggio streaming & \textbf{Streaming con $n$ noto ($m \ge 1$)}: seleziona $i_j$ con soglia condizionata dinamica $\frac{m - s}{n - j + 1}$ ($s$ elementi gi\`a estratti). Genera il campione nell'ordine naturale senza rimpiazzo n\'e sort finale; terminazione anticipata quando $s = m$. \\
\midrule
Priority Sampling (Random Keys) & \ref{chap:random-sampling-data-streams} & $\mathcal{O}(n \log m)$ tempo, $n$ chiamate RNG & $\mathcal{O}(m)$ (Min-Heap) & $1$ passaggio streaming & \textbf{Adatto sia con $n$ noto che incognito}: associa a ciascun elemento una chiave pseudo-casuale $u_i \sim \operatorname{Uniform}(0, 1)$ e trattiene in un Min-Heap di capienza $m$ i soli record a chiave massima. \\
\midrule
Reservoir Sampling (Algoritmo R) & \ref{chap:random-sampling-data-streams} & $\mathcal{O}(n)$ tempo, $n - m$ chiamate RNG & $\mathcal{O}(m)$ serbatoio & $1$ passaggio streaming & \textbf{Streaming con $n$ ignoto a priori ($m \ge 1$)}: accumula i primi $m$ elementi; per ogni $j > m$, include $i_j$ con probabilit\`a $m/j$ rimpiazzando un elemento estratto uniformemente nel buffer. \\
\midrule
Reservoir Sampling a Salti (Alg. Z) & \ref{chap:random-sampling-data-streams} & $\mathcal{O}(m(1 + \log(n/m)))$ tempo e RNG & $\mathcal{O}(m)$ serbatoio & $1$ passaggio streaming & \textbf{Streaming su big data ($n$ ignoto)}: modella analiticamente il numero di elementi da scartare prima del prossimo inserimento come variabile geometrica (\textit{skip distance}), saltando interi blocchi senza interrogare l'RNG. \\
\bottomrule
\end{tabular}%
}
\caption{Quadro comparativo delle strutture dati per dizionari, campionamento streaming e alberi cartesiani.}
\label{tab:sinottico-dizionari-streaming}
\end{table}

\FloatBarrier
\section{Intersezione di Posting Lists e Algoritmi di Partizionamento}
\label{sec:sinottico-intersezione-liste}

La Tabella~\ref{tab:sinottico-intersezione-liste} sintetizza le metodologie di intersezione di posting lists ordinate analizzate nel Capitolo~\ref{chap:posting-lists-intersection}, confrontando le tecniche sequenziali, le ricerche mirate e gli algoritmi basati sul partizionamento dello spazio di ricerca (Mutual Partition e approccio a blocchi a due livelli).

\begin{table}[H]
\centering
\small
\resizebox{\linewidth}{!}{%
\begin{tabular}{>{\raggedright\arraybackslash\bfseries}p{4.2cm} c >{\raggedright\arraybackslash}p{3.2cm} >{\raggedright\arraybackslash}p{2.5cm} >{\raggedright\arraybackslash}p{3.4cm} >{\raggedright\arraybackslash}p{6.6cm}}
\toprule
\textbf{Algoritmo / Tecnica} & \textbf{Cap.} & \textbf{Tempo / Confronti} & \textbf{Spazio Ausiliario} & \textbf{Accesso Memoria / I/O} & \textbf{Note Operative e Trade-off} \\
\midrule
Linear Merge-Scan & \ref{chap:posting-lists-intersection} & $\mathcal{O}(n + m)$ & $\mathcal{O}(1)$ & Sequenziale streaming & Scansione simmetrica a due puntatori; ottima per liste bilanciate ($n \approx m$). Massima localit\`a di cache e prefetching hardware; subottimale per $m \ll n$. \\
\midrule
Binary Search Intersection & \ref{chap:posting-lists-intersection} & $\mathcal{O}(m \log n)$ & $\mathcal{O}(1)$ & Salti casuali su $A$ & Cerca ciascun elemento di $B$ in $A$ con ricerca dicotomica; vantaggiosa solo per liste fortemente sbilanciate ($m < n/\log n$). Penalizzata da frequenti cache miss. \\
\midrule
Mutual Partition (Baeza-Yates) & \ref{chap:posting-lists-intersection} & $\Theta\left(m \left(1 + \log \frac{n}{m}\right)\right)$ & $\mathcal{O}(\log m)$ stack & Non sequenziale (ricorsivo) & \textbf{Ottimo asintotico teorico (Divide et Impera)}: sceglie il mediano di $A$, lo cerca in $B$ e partiziona contemporaneamente entrambe le sequenze in $A_1 \cap B_1$ e $A_2 \cap B_2$. Eguaglia il lower bound informativo $\Omega(m \log(n/m))$. Parallelizzabile, ma sfavorito su disco per la dispersione dei salti. \\
\midrule
Doubling / Galloping Search & \ref{chap:posting-lists-intersection} & $\mathcal{O}\left(m \log \frac{n}{m}\right)$ & $\mathcal{O}(1)$ & Forward streaming & \textbf{Ottimo asintotico iterativo}: salti esponenziali $2^k$ dall'ultimo match seguiti da ricerca binaria nell'intervallo ristretto. Eguaglia la complessit\`a di Mutual Partition eliminando la ricorsione e preservando la localit\`a spaziale. \\
\midrule
2-Level Block Partitioning (Skip Pointers) & \ref{chap:posting-lists-intersection} & $\mathcal{O}\left(\frac{n}{L} + m \cdot L\right) \to \mathcal{O}(\sqrt{n \cdot m})$ & $\mathcal{O}(n/L)$ sentinelle & Sequenziale a blocchi & \textbf{Approccio a due livelli gerarchico}: $A$ \`e suddivisa in blocchi di ampiezza $L$. Livello 1 = indice dei massimi $A'$; Livello 2 = blocchi contigui. Instrada $B$ su $A'$ e scandisce solo i blocchi attivi ($B_k \neq \emptyset$), saltando integralmente quelli inattivi ($0$ I/O). Con $L^* = \sqrt{n/m}$ tempo $\mathcal{O}(\sqrt{nm})$; standard de facto nei motori di ricerca reali. \\
\bottomrule
\end{tabular}%
}
\caption{Quadro comparativo degli algoritmi di intersezione per posting lists e partizionamento.}
\label{tab:sinottico-intersezione-liste}
\end{table}

\FloatBarrier
\section{Algoritmi e Strutture Dati per Stringhe in RAM}
\label{sec:sinottico-stringhe-ram}

La Tabella~\ref{tab:sinottico-stringhe-ram} espone le tecniche di ordinamento, ricerca e le strutture ad albero per insiemi di stringhe (Capitoli da~\ref{chap:string-sorting-search} a~\ref{chap:trie-compatti-msd-radix-sort-autocompletamento}).

\begin{table}[H]
\centering
\small
\resizebox{\linewidth}{!}{%
\begin{tabular}{>{\raggedright\arraybackslash\bfseries}p{4.2cm} c >{\raggedright\arraybackslash}p{3.2cm} >{\raggedright\arraybackslash}p{2.5cm} >{\raggedright\arraybackslash}p{3.4cm} >{\raggedright\arraybackslash}p{6.6cm}}
\toprule
\textbf{Algoritmo / Struttura} & \textbf{Cap.} & \textbf{Tempo di Ricerca} & \textbf{Spazio Struttura} & \textbf{Costo Costruzione} & \textbf{Note Operative e Punti di Forza} \\
\midrule
Interpolation Search & \ref{chap:string-sorting-search} & $\mathcal{O}(\log \log n)$ medio / $\mathcal{O}(n)$ w.c. & $\mathcal{O}(1)$ ausiliario & Array ordinato & Stima la posizione tramite interpolazione lineare; con vincolo a bucket (Bounded) assicura worst-case $\mathcal{O}(\log n)$. \\
\midrule
Multi-Key Quicksort & \ref{chap:string-sorting-search} & $\mathcal{O}(D + n \log n)$ & $\mathcal{O}(\log n + \max |s_i|)$ & In-place & Quicksort ternario a 3 vie ($<, =, >$) applicato al singolo carattere; non ricalcola caratteri gi\`a concordi ($D = \sum \text{lcp}$). \\
\midrule
Ternary Search Tree (TST) & \ref{chap:strutture-dati-stringhe-tst-radix-sort} & $\mathcal{O}(|P| + \log \Sigma)$ & $\mathcal{O}(n \cdot \bar{L})$ (3 puntatori/nodo) & $\mathcal{O}(D + n \log n)$ & Fonde la flessibilit\`a dei BST con la navigazione a prefisso dei Trie; ideale con alfabeti grandi e dizionari sparsi in RAM. \\
\midrule
LSD Radix Sort & \ref{chap:strutture-dati-stringhe-tst-radix-sort} & $\mathcal{O}(L \cdot (n + \Sigma))$ & $\mathcal{O}(n + \Sigma)$ ausiliario & — & Ordina da destra a sinistra con conteggio stabile; eccellente per chiavi a lunghezza fissa $L$ (es. indirizzi IP, hash). \\
\midrule
MSD Radix Sort & \ref{chap:trie-compatti-msd-radix-sort-autocompletamento} & $\mathcal{O}(D)$ (tempo ottimo) & $\mathcal{O}(\Sigma \cdot \max |s_i|)$ & — & Ordina dal carattere pi\`u significativo scendendo a bucket; richiede potatura a QuickSort per sottoinsiemi piccoli. \\
\midrule
Compact Trie (Patricia) & \ref{chap:trie-compatti-msd-radix-sort-autocompletamento} & $\mathcal{O}(|P|)$ & $\le 2n - 1$ nodi interni & $\mathcal{O}(D)$ & Comprime i cammini unari salvando l'indice di salto; elimina tutti i nodi interni a grado 1, garantendo linearit\`a nei nodi. \\
\midrule
Augmented Trie (Top-1) & \ref{chap:trie-compatti-msd-radix-sort-autocompletamento} & $\mathcal{O}(|P|)$ prefisso + $\mathcal{O}(1)$ risposta & $\mathcal{O}(n)$ & $\mathcal{O}(D)$ & Memorizza $\text{maxFreq}(u)$ e $\text{bestLeaf}(u)$; risponde a query di autocompletamento pesate in tempo $\mathcal{O}(1)$ alla fine del prefisso. \\
\midrule
Treap 2D per Stringhe & \ref{chap:trie-compatti-msd-radix-sort-autocompletamento} & Sub-lineare su range $I_P$ & $\mathcal{O}(n)$ & $\mathcal{O}(n \log n)$ & Combina ordine lessicografico BST con quote Max-Heap; consente pruning anticipato dei sottoalberi con frequenza inferiore al massimo corrente. \\
\bottomrule
\end{tabular}%
}
\caption{Quadro comparativo per ordinamento di stringhe e alberi digitali.}
\label{tab:sinottico-stringhe-ram}
\end{table}

\FloatBarrier
\section{Indicizzazione Avanzata, Memoria Gerarchica e Sottostringhe}
\label{sec:sinottico-indicizzazione-sottostringhe}

La Tabella~\ref{tab:sinottico-indicizzazione-sottostringhe} raccoglie le strutture di indicizzazione esterna per stringhe e la ricerca di sottostringhe con Suffix Tree, Suffix Array e Array LCP (Capitoli da~\ref{chap:ricerca-stringhe-memoria-gerarchica} a~\ref{ch:kasai-lcp-dizionari-esterni-occorrenze-vicine}).

\begin{table}[H]
\centering
\small
\resizebox{\linewidth}{!}{%
\begin{tabular}{>{\raggedright\arraybackslash\bfseries}p{4.2cm} c >{\raggedright\arraybackslash}p{3.2cm} >{\raggedright\arraybackslash}p{2.5cm} >{\raggedright\arraybackslash}p{3.4cm} >{\raggedright\arraybackslash}p{6.6cm}}
\toprule
\textbf{Algoritmo / Tecnica} & \textbf{Cap.} & \textbf{Tempo di Ricerca} & \textbf{Spazio Struttura} & \textbf{I/O Disco (DAM)} & \textbf{Note Operative e Punti di Forza} \\
\midrule
Three-Phase Blind Search & \ref{chap:ricerca-stringhe-memoria-gerarchica} & $\mathcal{O}(|P|)$ & $\mathcal{O}(n)$ scheletro compresso & $\mathbf{1}$ \textbf{I/O deterministico} & \textbf{Discesa cieca in RAM $\to$ 1 I/O su disco $\to$ Risalita in RAM}: risolve ricerca esatta o predecessore con 1 solo accesso a blocco. \\
\midrule
Front-Coding su Blocchi & \ref{chap:ricerca-stringhe-memoria-gerarchica} & $\mathcal{O}(|P| + B)$ scansione & $20\%\text{--}30\%$ della taglia originale & $\mathcal{O}(1)$ blocco caricato & Comprime prefissi comuni consecutivi salvando $\langle \text{lcp}, \text{suffisso} \rangle$; record header non compresso per accesso casuale. \\
\midrule
Suffix Tree (Ukkonen) & \ref{ch:indicizzazione-ricerca-sottostringhe-suffix-array-tree} & $\mathcal{O}(|P|)$ o $\mathcal{O}(|P| \log \Sigma)$ & $20n\text{--}45n$ byte (ingombrante) & RAM-oriented & Albero compatto di tutti i suffissi; risolve occorrenze, sottostringhe ripetute e palindromi in tempo ottimo; footprint RAM elevato. \\
\midrule
Suffix Array (Dicotomico) & \ref{ch:indicizzazione-ricerca-sottostringhe-suffix-array-tree} & $\mathcal{O}(|P| \log n)$ & $n \log_2 n$ bit ($4n\text{--}8n$ byte) & RAM-oriented & Vettore compatto dei puntatori ai suffissi ordinati; individua l'intervallo contiguo $[L, R]$ con due ricerche binarie indipendenti. \\
\midrule
Suffix Array + LCP & \ref{ch:indicizzazione-ricerca-sottostringhe-suffix-array-tree} & $\mathcal{O}(|P| + \log n)$ & $SA + LCP \approx 8n$ byte & RAM-oriented & Sfrutta i prefissi condivisi memorizzati nell'array LCP per non riesaminare i caratteri certificati durante il dimezzamento binario. \\
\midrule
Algoritmo di Kasai et al. & \ref{ch:kasai-lcp-dizionari-esterni-occorrenze-vicine} & Costruzione: $\Theta(n)$ ($\le 3n$ confronti) & $\mathcal{O}(n)$ ausiliario ($Rank[1\dots n]$) & $0$ I/O (in RAM) & Calcola l'intero array LCP in tempo deterministico lineare processando i suffissi nell'ordine posizionale naturale $i = 1 \dots n$. \\
\midrule
Closest Occurrences (Tagged Merge) & \ref{ch:kasai-lcp-dizionari-esterni-occorrenze-vicine} & $\mathcal{O}(\text{range} + occ \log occ)$ & $\mathcal{O}(occ_P + occ_Q)$ & $0$ I/O (in RAM) & Localizza su $SA$ gli intervalli con sentinelle per $P$ e $Q$; fonde le liste e trova la minima distanza con scansione lineare adiacente. \\
\midrule
Distinct Strings Count (Blind Trie) & \ref{ch:kasai-lcp-dizionari-esterni-occorrenze-vicine} & $\mathcal{O}(N_{\text{tot}})$ & $\mathcal{O}(d)$ nodi in RAM & $d$ I/O ($1$ per stringa distinta) & Conteggio stringhe distinte in memoria esterna: mantiene il solo scheletro compresso in RAM, verificando su disco solo i nuovi candidati. \\
\bottomrule
\end{tabular}%
}
\caption{Quadro comparativo per indicizzazione gerarchica, strutture per sottostringhe e algoritmi su Suffix Array.}
\label{tab:sinottico-indicizzazione-sottostringhe}
\end{table}
